% NOETHER PAPER 4 — KOREAN TRANSLATION-PRODUCER DRAFT — UNCHECKED
% Unit: T04-U17; current-authority whole lines 3889--3908
% Source slice: 1,540 LF UTF-8 bytes; SHA-256 567E06533B6D2EBD77AA6E15726ABFC73135B2AC5C93E9EE3C46B5766387BA43
% Pointer: NOETH-DE-AUTH-v004-20260804; SHA-256 A1C62FDACAA34DFC1B806DC18258F2E732539F7AE9D85AA4BD9E1067B8749D9F
% This is producer translation only: no source/Korean/formula review, compilation, or rendering.

\section*{\S\ 4. 분해 항등식.}
이제 임의의 기호열과 변수열 사이의 항등식 가운데 가장 일반적인 경우를 생각한다. 이 항등식의 종결식에는 대입 (11.)을 하지 않은 채로 하나의 열 $p_\tau$가 그 개별 열 $y^{(1)}\cdots y^{(\tau)}$로 분해되어 나타난다. 우리는 이러한 항등식을 ``분해 항등식''이라 부른다. 앞 절 끝의 주석에 이어서, 먼저 기호열이 두 개뿐인 경우에 이 항등식을 구한다. 즉, 식 $\pair{q_\rho A^\sigma}{B_{\rho+\sigma-\tau}p_\tau}$을 전개한다. 이때 다음과 같이 놓는다.
\begin{equation}
\begin{aligned}
q_{\rho-\alpha}^{(i_1\ldots i_\alpha)}&\sim p_{n-\rho}\,y^{(i_1)}\cdots y^{(i_\alpha)},\qquad
p_{\tau-\alpha}^{(i_1\ldots i_\alpha)}=y^{(i_{\alpha+1})}\cdots y^{(i_\tau)},\qquad
p_\tau=y^{(1)}\cdots y^{(\tau)}.
\end{aligned}
\tag{22}
\end{equation}
수 $\rho,\sigma,\tau$에 관해서는 네 경우를 구별해야 한다.
\[
\begin{array}{llll}
1)&\rho+\sigma\leqq n,& \tau\leqq\rho,& \tau\leqq\sigma;\\[2pt]
2)&\rho+\sigma\leqq n,& \tau\geqq\rho,& \tau\leqq\sigma;\\[2pt]
3)&\rho+\sigma\leqq n,& \tau\leqq\rho,& \tau>\sigma;\\[2pt]
4)&\rho+\sigma\leqq n,& \tau\geqq\rho,& \tau\geqq\sigma,
\end{array}
\]
이 가운데 첫째 경우를 택하여 다룬 뒤 나머지 경우를 짧게 특징짓겠다. 열 $y^{(1)}y^{(2)}\cdots y^{(\tau)}B_{\rho+\sigma-\tau}$를 $y^{(1)}$과 $y^{(2)}\cdots y^{(\tau)}B_{\rho+\sigma-\tau}$로 분해하면, (20.), (10.), (9.)에 따라 다음이 나온다.
